\subsection{Scenario Validation}\label{sec:scenario-validation}

We construct five synthetic scenarios with known diversity structure:
\begin{enumerate}[nosep]
    \item \textbf{Pure noise}: Random word salad with no learnable structure.
    \item \textbf{Multi-incoherent}: Responses from 5 distinct modes, each internally incoherent (scrambled words within a template).
    \item \textbf{Multi-mode}: Responses from 5 distinct modes, each internally coherent (e.g., recipe, poem, code).
    \item \textbf{One-mode}: All responses from a single coherent mode (paraphrases of the same content).
    \item \textbf{Mixed}: A mixture of coherent and incoherent responses.
\end{enumerate}

Each scenario contains 10 responses per prompt, 5 prompts per scenario, with 100 permutations and seed 42.

\begin{table}[t]
\centering
\caption{Scenario validation metrics (mean across 5 prompts, 100 permutations, $n=10$ responses, per-byte).
We report several candidate scalars: unconditional surprise $a_1$, the asymptotic floor $a_n$, coherence $C$, our recommended score $D_{Ca_n} = C \times a_n$ (\textbf{bold}), the excess entropy $E$ (Appendix~\ref{app:excess-entropy}), $C \times E$, and the coherence spread $\sigma_\ell$. $D_{Ca_n}$ correctly ranks multi-mode above one-mode and suppresses incoherent scenarios via $C$.
Both $E$ and $C \times E$ go negative for the mixed scenario on GPT-2.}
\label{tab:scenarios}
\small
% Generated by: scripts/generate_scenario_table.py
% Data sources: results/scenario_metrics_v2_100perm.json
%               results/scenario_metrics_v2_qwen_100perm.json
\begin{tabular}{@{}l rrrrrrr rrrrrrr@{}}
\toprule
& \multicolumn{7}{c}{\textbf{GPT-2 (124M)}} & \multicolumn{7}{c}{\textbf{Qwen2.5-32B}} \\
\cmidrule(lr){2-8} \cmidrule(lr){9-15}
Scenario & $a_1$ & $a_n$ & $C$ & $\mathbf{D_{a_\infty}}$ & $E$ & $C\!\times\!E$ & $\sigma_\ell$ & $a_1$ & $a_n$ & $C$ & $\mathbf{D_{a_\infty}}$ & $E$ & $C\!\times\!E$ & $\sigma_\ell$ \\
\midrule
Pure noise & 7.32 & 6.84 & 0.006 & \textbf{0.04} & 0.55 & 0.003 & 0.20 & 6.85 & 6.56 & 0.009 & \textbf{0.06} & 0.44 & 0.004 & 0.09 \\
Multi-incoher. & 2.99 & 1.66 & 0.12 & \textbf{0.20} & 5.64 & 0.68 & 0.51 & 1.93 & 1.19 & 0.27 & \textbf{0.33} & 2.42 & 0.66 & 0.84 \\
Multi-mode & 1.02 & 0.53 & 0.50 & \textbf{0.26} & 1.55 & 0.75 & 0.08 & 0.71 & 0.26 & 0.62 & \textbf{0.16} & 1.21 & 0.74 & 0.09 \\
One-mode & 1.02 & 0.22 & 0.50 & \textbf{0.11} & 1.24 & 0.61 & 0.06 & 0.74 & 0.15 & 0.60 & \textbf{0.09} & 0.85 & 0.51 & 0.08 \\
Mixed & 1.62 & 1.71 & 0.29 & \textbf{0.49} & $-$1.27 & $-$0.34 & 1.12 & 0.77 & 0.56 & 0.57 & \textbf{0.32} & 0.69 & 0.39 & 0.47 \\
\bottomrule
\end{tabular}

\end{table}

\paragraph{Results.} Table~\ref{tab:scenarios} summarizes the metrics.
Figure~\ref{fig:scenario-curves} shows the $a_k$ curves for all scenarios.

\begin{figure}[t]
\centering
\includegraphics[width=\textwidth,height=0.4\textheight,keepaspectratio]{comparison_v3/comparison_ak_curves_all.png}
\caption{Progressive conditional surprise curves ($a_k$, per-byte) for all five scenarios, comparing GPT-2 (top row) and Qwen2.5-3B (bottom row).
In each panel, faint colored curves are individual permutations (100 per prompt), medium colored curves are per-prompt averages, and the bold black curve is the mean across prompts.
Pure noise curves are flat (no learnable structure); multi-mode curves show progressive decline; one-mode curves decline steeply then flatten.}
\label{fig:scenario-curves}
\end{figure}

The coherence term $C$ correctly separates coherent from incoherent scenarios: $C \approx 0.5$--$0.6$ for multi-mode and one-mode, $C \approx 0.1$--$0.3$ for multi-incoherent, and $C \approx 0.01$ for pure noise.
The coherence spread $\sigma_\ell$ discriminates mixed from pure scenarios: $\sigma_\ell > 1.0$ for mixed (GPT-2), reflecting the within-set heterogeneity between coherent and incoherent responses.

$C$ consistently improves with model strength (Qwen2.5-3B assigns higher coherence to genuinely coherent text than GPT-2). $D_{Ca_n}$ correctly ranks multi-mode above one-mode on both models, and suppresses pure noise and multi-incoherent via the $C$ factor.

